% sec1_1_part3.tex — Section 1.1 continued: Classical Regression Model for
% Random Samples; "Fixed" Regressors (book pp. 12–14)

Since the $(i, j)$ element of the $n \times n$ matrix
$\bm{\varepsilon}\bm{\varepsilon}'$ is $\varepsilon_i \varepsilon_j$,
Assumption~1.4 can be written compactly as
\begin{equation}\label{eq:1.1.14}
  \E(\bm{\varepsilon}\bm{\varepsilon}' \mid \vec{X})
    = \sigma^{2} \vec{I}_{n}.
\end{equation}
The discussion of the previous paragraph shows that the assumption can also be
written as
\[
  \Var(\bm{\varepsilon} \mid \vec{X}) = \sigma^{2} \vec{I}_{n}.
\]
However, \eqref{eq:1.1.14} is the preferred expression, because the more
convenient measure of variability is second moments (such as
$\E(\varepsilon_i^{2} \mid \vec{X})$) rather than variances.  This point will
become clearer when we deal with the large sample theory in the next chapter.
Assumption~1.4 is sometimes called the \textbf{spherical} error variance
assumption because the $n \times n$ matrix of second moments (which are also
variances and covariances) is proportional to the identity matrix
$\vec{I}_{n}$.  This assumption will be relaxed later in this chapter.

%% ---------------------------------------------------------------
\subsection*{The Classical Regression Model for Random Samples}

The sample $(\vec{y}, \vec{X})$ is a \textbf{random sample} if
$\{y_i, \vec{x}_i\}$ is i.i.d.\ (independently and identically distributed)
across observations.  Since by Assumption~1.1 $\varepsilon_i$ is a function of
$(y_i, \vec{x}_i)$ and since $(y_i, \vec{x}_i)$ is independent of
$(y_j, \vec{x}_j)$ for $j \neq i$,
$(\varepsilon_i, \vec{x}_i)$ is independent of $\vec{x}_j$ for $j \neq i$.
So
\begin{align}
  \E(\varepsilon_i \mid \vec{X})
    &= \E(\varepsilon_i \mid \vec{x}_i), \notag \\
  \E(\varepsilon_i^{2} \mid \vec{X})
    &= \E(\varepsilon_i^{2} \mid \vec{x}_i), \notag \\
  \text{and}\quad
  \E(\varepsilon_i \varepsilon_j \mid \vec{X})
    &= \E(\varepsilon_i \mid \vec{x}_i)\,
       \E(\varepsilon_j \mid \vec{x}_j)
       \quad (\text{for } i \neq j).
  \label{eq:1.1.15}
\end{align}
(Proving the last equality in \eqref{eq:1.1.15} is a review question.)
Therefore, Assumptions~1.2 and 1.4 reduce to
\begin{align}
  \text{Assumption 1.2:}\quad
    & \E(\varepsilon_i \mid \vec{x}_i) = 0
      \quad (i = 1, 2, \ldots, n),
    \label{eq:1.1.16} \\
  \text{Assumption 1.4:}\quad
    & \E(\varepsilon_i^{2} \mid \vec{x}_i) = \sigma^{2} > 0
      \quad (i = 1, 2, \ldots, n).
    \label{eq:1.1.17}
\end{align}

The implication of the identical distribution aspect of a random sample is that
the joint distribution of $(\varepsilon_i, \vec{x}_i)$ does not depend on $i$.
So the \emph{un}conditional second moment $\E(\varepsilon_i^{2})$ is constant
across $i$ (this is referred to as \textbf{unconditional homoskedasticity})
and the functional form of the conditional second moment
$\E(\varepsilon_i^{2} \mid \vec{x}_i)$ is the same across $i$.  However,
Assumption~1.4---that the \emph{value} of the conditional second moment is the
same across $i$---does not follow.  Therefore, Assumption~1.4 remains
restrictive for the case of a random sample; without it, the conditional second
moment $\E(\varepsilon_i^{2} \mid \vec{x}_i)$ can differ across $i$ through
its possible dependence on $\vec{x}_i$.  To emphasize the distinction, the
restrictions on the conditional second moments, \eqref{eq:1.1.12} and
\eqref{eq:1.1.17}, are referred to as \textbf{conditional homoskedasticity}.

%% ---------------------------------------------------------------
\subsection*{``Fixed'' Regressors}

We have presented the classical linear regression model, treating the
regressors as random.  This is in contrast to the treatment in most textbooks,
where $\vec{X}$ is assumed to be ``fixed'' or deterministic.  If $\vec{X}$ is
fixed, then there is no need to distinguish between the conditional
distribution of the error term, $f(\varepsilon_i \mid \vec{x}_1, \ldots,
\vec{x}_n)$, and the unconditional distribution, $f(\varepsilon_i)$, so that
Assumptions~1.2 and 1.4 can be written as
\begin{align}
  \text{Assumption 1.2:}\quad
    & \E(\varepsilon_i) = 0
      \quad (i = 1, \ldots, n),
    \label{eq:1.1.18} \\
  \text{Assumption 1.4:}\quad
    & \E(\varepsilon_i^{2}) = \sigma^{2}
      \quad (i = 1, \ldots, n); \notag \\
    & \E(\varepsilon_i \varepsilon_j) = 0
      \quad (i, j = 1, \ldots, n;\; i \neq j).
    \label{eq:1.1.19}
\end{align}

Although it is clearly inappropriate for a nonexperimental science like
econometrics, the assumption of fixed regressors remains popular because the
regression model with fixed $\vec{X}$ can be interpreted as a set of statements
conditional on $\vec{X}$, allowing us to dispense with ``$\mid\;\vec{X}$''
from the statements such as Assumptions~1.2 and 1.4 of the model.

However, the economy in the notation comes at a price.  It is very easy to miss
the point that the error term is being assumed to be uncorrelated with current,
past, and future regressors.  Also, the distinction between the unconditional
and conditional homoskedasticity gets lost if the regressors are deterministic.
Throughout this book, the regressors are treated as random, and, unless
otherwise noted, statements conditional on $\vec{X}$ are made explicit by
inserting ``$\mid\;\vec{X}$.''
